\documentclass[11pt]{article}
\usepackage[T1]{fontenc}
\usepackage{lmodern}
\usepackage{amsmath,amssymb,amsthm}
\usepackage[margin=1in]{geometry}
\usepackage[hidelinks]{hyperref}
\usepackage{microtype}

\newtheorem{theorem}{Theorem}
\newcommand{\spec}{\operatorname{spec}}

\title{Two isospectral periodic Jacobi matrices\texorpdfstring{\\}{ }%
with different diagonals}
\author{Feiyu (Independent researcher)}
\date{}

\begin{document}
\maketitle

\begin{abstract}
Two bilateral Jacobi matrices of least period two both have spectrum
$[-5,-3]\cup[3,5]$. The diagonal vanishes everywhere for one matrix and
nowhere for the other. Both bands have positive length, so the pair
disproves the mirror-uniqueness clause of Conjecture 00000008581 in the
absence of degenerate bands.
\end{abstract}

\section{The conjecture as stated}

The English version of Conjecture 00000008581 reads as follows. The
Chinese version appears in the same \href{https://github.com/The-Last-Math-Competition/The-Last-Math-Competition/blob/8c9cc2fb29ed8a59d7daa18ed6ba387a77cf897f/conjectures/00000008581.md}{official statement}.

\begin{quote}
Definition: Finite-band periodic Jacobi matrices: real symmetric
tridiagonal matrices whose spectrum is a union of finitely many
intervals. Conjecture: The criterion for N-band potentials is given by
the Sodin-Yuditskii transfer function condition (an S-Y criterion law);
the solution of the inverse spectral problem is unique up to mirror
symmetry (a mirror uniqueness law); uniqueness fails in the hole case
(spectral intervals degenerating to points, a hole exception law); and
the parameter dimension of hole examples is 1 (a one-dimensional
parameter law). (finite-band Jacobi Sodin-Yuditskii mirror)
\end{quote}

\section{Reading of the statement}

A bilateral Jacobi matrix acts on $\ell^2(\mathbb Z;\mathbb C)$ by
\[
 (Ju)_n = a_{n-1}u_{n-1}+b_nu_n+a_nu_{n+1},
 \qquad a_n>0,\quad b_n\in\mathbb R.
\]
The coefficient sequences are periodic with a common period. Such
coefficients are bounded, so the displayed formula gives a bounded
operator. Since its entries are real and symmetric, it is self-adjoint
by symmetry on the dense subspace of finitely supported sequences.
We use the operator spectrum, the set of complex $z$ such that
$zI-J$ has no bounded inverse. This is the bilateral inverse problem
for periodic Jacobi matrices, as in \cite[\S 1 and \S 2.14]{dks}.

Both versions assert uniqueness up to mirror symmetry and mention an
exception when an interval degenerates to a point. In particular, for
any two such period-two matrices and any $\alpha<\beta<\gamma<\delta$,
if their spectra both equal
$[\alpha,\beta]\cup[\gamma,\delta]$, their entries must agree after a
reflection of the index line, allowing a translation of its origin
as well. The counterexample disproves even this restricted claim.

\clearpage
\section{Counterexample}

\begin{theorem}
The period-two Jacobi matrices with coefficients
\[
\begin{array}{c|cccc}
 & a_{2k} & a_{2k+1} & b_{2k} & b_{2k+1} \\ \hline
 J_1 & 1 & 4 & 0 & 0 \\
 J_2 & 2 & 2 & 3 & -3
\end{array}
\qquad (k\in\mathbb Z)
\]
have common spectrum $[-5,-3]\cup[3,5]$. No simultaneous relabeling
of rows and columns identifies their entries.
\end{theorem}

\begin{proof}
Write the coefficients of a matrix of this form as $(a,c,b,-b)$.
Group the even and odd sites, putting $x_k=u_{2k}$ and $y_k=u_{2k+1}$.
This identifies the original space with two copies of
$\ell^2(\mathbb Z;\mathbb C)$. On each copy let $S$ be the shift
$(Sx)_k=x_{k+1}$, with inverse $(S^{-1}x)_k=x_{k-1}$.
In the even and odd coordinates the operator and its square are
\[
 J=\begin{pmatrix}bI&aI+cS^{-1}\\aI+cS&-bI\end{pmatrix},
 \qquad
 J^2=(a^2+c^2+b^2)I+ac\,(T\oplus T),\qquad T=S+S^{-1}.
\]
Here $I$ denotes the identity on the relevant space. For either row
in the table, $a^2+c^2+b^2=17$ and $ac=4$.

We first compute the spectrum of $S$. Since $S$ and $S^{-1}$ are
isometries, both spectra lie in the closed unit disk. If
$z\in\spec(S)$, then $z\ne 0$ and $z^{-1}\in\spec(S^{-1})$.
It follows that $|z|=1$. Conversely, for $|z|=1$, a solution of
$(zI-S)u=\delta_0$ would satisfy
$|u_{n+1}|=|u_n|$ for $n\ne0$. The constant modulus on each of the
two half-lines must vanish by square summability. This gives
$u_0=u_1=0$, contradicting $z u_0-u_1=1$. Thus $\spec(S)$ is the
entire unit circle.

For any $t\in\mathbb C$, let $z,w$ be the roots of $X^2-tX+1$.
Then
\[
 tI-T=-S^{-1}(S-zI)(S-wI).
\]
The two factors on the right commute, so their product is invertible
exactly when both are invertible. Consequently $t\in\spec(T)$
exactly when $t=z+z^{-1}$ for some $|z|=1$. These values fill
$[-2,2]$. The spectrum of $T\oplus T$ equals that of $T$ because its
resolvent is invertible exactly when the common block resolvent is.
It follows that $\spec(J^2)=17+4[-2,2]=[9,25]$ for both matrices.

Let $(Rx)_k=x_{-k}$ on one copy. The invertible map
$(x,y)\mapsto(Ry,-Rx)$ conjugates $J$ to $-J$, as follows from
$RS=S^{-1}R$ and the displayed block matrix. Thus $\spec(J)$ is
symmetric under negation. Polynomial spectral mapping for the
square then gives
\[
 z\in\spec(J)\quad\Longleftrightarrow\quad z^2\in[9,25],
\]
which yields precisely the two claimed bands. Both have length $2$,
and their intervening gap is $(-3,3)$.

Period one is impossible for $J_1$ because the bonds alternate and
for $J_2$ because the diagonal alternates. A simultaneous
relabeling preserves the set of diagonal values, which is $\{0\}$
for $J_1$ and $\{-3,3\}$ for $J_2$. This rules out every reflection
and translation of the index line.
\end{proof}

\clearpage
\section{Lean formalization}

\begingroup\raggedright
The Lean project uses Mathlib's $\ell^2$ type \texttt{lp},
bounded operators \texttt{ContinuousLinearMap}, and
\texttt{spectrum} over $\mathbb C$. The structure \texttt{Jacobi}
represents the period-two subclass with alternating diagonal
$b,-b$ and includes strict positivity of both bonds.

\texttt{Jacobi.period\_two} certifies periodicity.
\texttt{Jacobi.siteOperator\_apply} verifies the tridiagonal action
at every integer site. The theorems \texttt{Jacobi.entry\_symmetric}
and \texttt{Jacobi.entry\_eq\_operator} verify symmetry and relate
the entries to the operator on coordinate vectors. The proof constructs
a continuous linear equivalence between the scalar sequence space and
its even and odd components, then transfers the spectrum across
this equivalence.

The theorems \texttt{first\_spectrum} and \texttt{second\_spectrum}
identify both spectra exactly, including the band endpoints.
\texttt{no\_site\_relabeling} rules out even an arbitrary map
of site indices inducing equality of the two entry arrays.
\texttt{not\_mirrorUniqueness} negates the universally quantified
two-band consequence stated above.

The entry point is \texttt{import Submission}. The theorem
\texttt{conjecture\_00000008581\_false} negates the conjunction with
any other clauses. Lean and Mathlib are pinned to v4.33.1. The
theorem's only axiom dependencies are \texttt{propext},
\texttt{Classical.choice}, and \texttt{Quot.sound}.

\par
\endgroup
\begin{thebibliography}{1}
\bibitem{dks}
D.~Damanik, R.~Killip, and B.~Simon,
\emph{Perturbations of orthogonal polynomials with periodic recursion
coefficients}, Ann. of Math. (2) \textbf{171} (2010), 1931--2010.
\href{https://annals.math.princeton.edu/wp-content/uploads/annals-v171-n3-p11-p.pdf}{Published article}.
\end{thebibliography}

\end{document}
